\chapter*{Introduction}
\addcontentsline{toc}{chapter}{Introduction}

The introduction is a mandatory part of a bachelor's/master's thesis. The introduction is an introduction to the topic. It elaborates on the chosen topic, briefly places it in context (this may include a description of the motivation for writing the thesis), and answers the question of why the topic was chosen. It places the topic in context and justifies its necessity and the timeliness of addressing it. It contains an explicitly stated aim of the thesis. The text of the aim must match the text given in the thesis assignment, i.e. the text given in the SIS system, which is also given in the Abstract section.

The introduction also includes a brief presentation of how the thesis was processed (the processing method is covered in detail in a separate part of the main text). The introduction may also include a description of the motivation for writing the thesis.

The introduction to a master's thesis must be more elaborate -- see your faculty's thesis guidelines for details.

Several sample chapters follow, recommending how a final thesis should be typeset. They primarily describe the use of the \TeX{} template, but the general advice will also serve users of other systems well.
